\documentclass[12pt,a4paper]{article}
\usepackage[utf8]{inputenc}
\usepackage[T1]{fontenc}
\usepackage{amsmath,amssymb,amsthm}
\usepackage{physics}
\usepackage{geometry}
\usepackage{enumitem}
\usepackage{hyperref}
\usepackage{fancyhdr}
\usepackage{titlesec}

\geometry{margin=1in}
\pagestyle{fancy}
\fancyhf{}
\rhead{Lecture 4}
\lhead{Advanced Quantum Mechanics}
\cfoot{\thepage}

\titleformat{\section}{\large\bfseries}{}{0em}{}
\titleformat{\subsection}{\normalsize\bfseries}{}{0em}{}

\newtheorem{definition}{Definition}

\title{\textbf{Lecture 4: The Harmonic Oscillator and Spin} \\ \large Wave Functions, Creation/Annihilation Operators, and Spin-$\frac{1}{2}$}
\author{Based on Leonard Susskind's \textit{Advanced Quantum Mechanics}}
\date{}

\begin{document}
\maketitle
\thispagestyle{fancy}

% ============================================================
\section{1. Algebraic Structure of the Harmonic Oscillator}
% ============================================================

The harmonic oscillator Hamiltonian is:
\[
    H = \frac{p^2}{2m} + \frac{1}{2}m\omega^2 x^2.
\]
The entire spectrum can be derived algebraically using the creation and annihilation operators.

\begin{definition}[Creation and annihilation operators]
Define:
\[
    a = \sqrt{\frac{m\omega}{2\hbar}}\left(x + \frac{ip}{m\omega}\right), \qquad
    a^\dagger = \sqrt{\frac{m\omega}{2\hbar}}\left(x - \frac{ip}{m\omega}\right).
\]
These satisfy the canonical commutation relation:
\[
    [a, a^\dagger] = 1.
\]
\end{definition}

\subsection{1.1 The number operator}

The \textbf{number operator} is:
\[
    N = a^\dagger a, \qquad N\ket{n} = n\ket{n}.
\]
The Hamiltonian expressed in terms of $N$:
\[
    H = \hbar\omega\left(N + \frac{1}{2}\right) = \hbar\omega\left(a^\dagger a + \frac{1}{2}\right).
\]

\subsection{1.2 Commutation relations with $N$}

\begin{align}
    [N, a] &= -a, \\
    [N, a^\dagger] &= +a^\dagger.
\end{align}
These show that $a$ lowers the eigenvalue of $N$ by one, and $a^\dagger$ raises it by one:
\[
    a^\dagger\ket{n} = \sqrt{n+1}\,\ket{n+1}, \qquad a\ket{n} = \sqrt{n}\,\ket{n-1}.
\]

\subsection{1.3 The ground state}

The ground state $\ket{0}$ is defined by $a\ket{0} = 0$.  All excited states are built by repeated application of $a^\dagger$:
\[
    \ket{n} = \frac{(a^\dagger)^n}{\sqrt{n!}}\ket{0}.
\]

% ============================================================
\section{2. Harmonic Oscillator Wave Functions}
% ============================================================

\subsection{2.1 Ground-state wave function}

The condition $a\ket{0} = 0$ translates in position space to:
\[
    \left(x + \frac{\hbar}{m\omega}\frac{d}{dx}\right)\psi_0(x) = 0.
\]
The solution is a Gaussian:
\[
    \boxed{\psi_0(x) = \left(\frac{m\omega}{\pi\hbar}\right)^{1/4}\exp\left(-\frac{m\omega x^2}{2\hbar}\right).}
\]

\subsection{2.2 Excited-state wave functions}

Higher wave functions are obtained by applying the creation operator:
\[
    \psi_n(x) = \frac{1}{\sqrt{n!}}\left(\frac{m\omega}{2\hbar}\right)^{n/2}\left(x - \frac{\hbar}{m\omega}\frac{d}{dx}\right)^n \psi_0(x).
\]
These are Gaussian functions multiplied by Hermite polynomials $H_n$:
\[
    \psi_n(x) \propto H_n\!\left(\sqrt{\frac{m\omega}{\hbar}}\,x\right)\exp\left(-\frac{m\omega x^2}{2\hbar}\right).
\]
The $n$-th wave function has exactly $n$ nodes, consistent with the general principle that higher-energy states oscillate more rapidly.

% ============================================================
\section{3. Ubiquity of the Harmonic Oscillator}
% ============================================================

The harmonic oscillator appears throughout physics:
\begin{itemize}[nosep]
    \item \textbf{Molecular vibrations}: atoms in a molecule oscillate about equilibrium positions.
    \item \textbf{Phonons}: quantized vibrations of a crystal lattice.
    \item \textbf{Photons}: each mode of the electromagnetic field is a harmonic oscillator; the quanta are photons.
    \item \textbf{Any stable equilibrium}: expanding any smooth potential about a minimum gives $V(x) \approx V_0 + \tfrac{1}{2}V''(x_0)(x - x_0)^2$, which is a harmonic oscillator.
\end{itemize}

% ============================================================
\section{4. Spin-$\frac{1}{2}$}
% ============================================================

\subsection{4.1 Intrinsic angular momentum}

Particles possess an intrinsic angular momentum called \textbf{spin}, which has no classical analog of a spinning object.  For an electron (spin-$\frac{1}{2}$), the spin operators satisfy:
\[
    [S_i, S_j] = i\hbar\,\epsilon_{ijk}\,S_k,
\]
the same algebra as orbital angular momentum.  The spin quantum number is $s = \frac{1}{2}$, giving two states:
\[
    S_z\ket{\uparrow} = +\frac{\hbar}{2}\ket{\uparrow}, \qquad S_z\ket{\downarrow} = -\frac{\hbar}{2}\ket{\downarrow}.
\]

\subsection{4.2 Pauli matrices}

The spin operators for spin-$\frac{1}{2}$ are $S_i = \frac{\hbar}{2}\sigma_i$, where the \textbf{Pauli matrices} are:
\[
    \sigma_x = \begin{pmatrix} 0 & 1 \\ 1 & 0 \end{pmatrix}, \qquad
    \sigma_y = \begin{pmatrix} 0 & -i \\ i & 0 \end{pmatrix}, \qquad
    \sigma_z = \begin{pmatrix} 1 & 0 \\ 0 & -1 \end{pmatrix}.
\]
Key properties:
\begin{itemize}[nosep]
    \item $\sigma_i^2 = I$ for each $i$.
    \item $\sigma_i\sigma_j + \sigma_j\sigma_i = 2\delta_{ij}I$ (anti-commutation).
    \item $\sigma_i\sigma_j - \sigma_j\sigma_i = 2i\epsilon_{ijk}\sigma_k$ (commutation).
\end{itemize}

% ============================================================
\section{5. Bosons and Fermions --- Preview}
% ============================================================

Particles come in two fundamental types:
\begin{itemize}[nosep]
    \item \textbf{Bosons}: integer spin ($s = 0, 1, 2, \ldots$).  Examples: photons, phonons, mesons.
    \item \textbf{Fermions}: half-integer spin ($s = \frac{1}{2}, \frac{3}{2}, \ldots$).  Examples: electrons, protons, neutrons.
\end{itemize}
The distinction is \emph{not} merely a label---it has profound consequences for the behavior of identical particles, which will be the subject of the next lecture.

\medskip
\noindent\textit{The spin-statistics connection---that integer-spin particles are bosons and half-integer-spin particles are fermions---is one of the deepest results in theoretical physics.  Its consequences include the Pauli exclusion principle and the structure of the periodic table.}

\end{document}
